% GENERATED by scripts/gen_blueprint.py — do not edit.
\chapter{General tournament theory}
\label{bp:ch_general}

\section{The substitution lower bound}

\begin{theorem}[\code{omegabar\_lexprod\_ge}]
  \label{res:omegabar_lexprod_ge}
  \uses{def:arc, def:backedge, def:lexprod, def:ltp, def:omegabar, def:tournament}
  \rocq{Digraph.invariants_advanced.substitution.omegabar_lexprod_ge}
  \rocqok

\begin{verbatim}
Theorem omegabar_lexprod_ge :
  ω̄(S) + ω̄(H) <= ω̄(lexprod_tournament S H) + 1.
\end{verbatim}
\href{https://github.com/LLM4Rocq/digraph-theory/blob/main/theories/invariants_advanced/substitution.v\#L168}{Source: \code{theories/invariants\_advanced/substitution.v:168}}


\paragraph{Decoded.} $\bar\omega(S) + \bar\omega(H) \le \bar\omega(S[H]) + 1$ for any
tournaments $S, H$ (quoted inside a section fixing nonempty
\code{S}, \code{H}): substitution composes criticality
levels. This is what makes the tower $\mathrm{AC}_n \to
\mathrm{AC}_n[C_3] \to \mathrm{AC}_n[\mathrm{AC}_n]$ work.

\end{theorem}

\section{Vertex-transitive tournaments: one deletion decides criticality}

\begin{theorem}[\code{vt\_kcritical}]
  \label{res:vt_kcritical}
  \uses{def:arc, def:aut, def:backedge, def:kcritical, def:ltp, def:omegabar, def:subdel, def:tournament, def:vt}
  \rocq{Digraph.invariants_advanced.transitive.vt_kcritical}
  \rocqok

\begin{verbatim}
Theorem vt_kcritical (k : nat) (v0 : T) :
  vertex_transitiveb (T : diGraphType) ->
  kcritical k T = (ω̄(T) == k) && (ω̄(del_tournament v0) == k.-1).
\end{verbatim}
\href{https://github.com/LLM4Rocq/digraph-theory/blob/main/theories/invariants_advanced/transitive.v\#L57}{Source: \code{theories/invariants\_advanced/transitive.v:57}}


\paragraph{Decoded.} In a vertex-transitive tournament all single-vertex deletions have
the same $\bar\omega$ (second statement), so $k$-criticality reduces
to the value of $\bar\omega(T)$ and of one deletion
$\bar\omega(T - v_0)$ (first statement, a Boolean equivalence).

\end{theorem}

\begin{theorem}[\code{omegabar\_del\_vt}]
  \label{res:omegabar_del_vt}
  \uses{def:arc, def:aut, def:backedge, def:ltp, def:omegabar, def:subdel, def:tournament, def:vt}
  \rocq{Digraph.invariants_advanced.transitive.omegabar_del_vt}
  \rocqok

\begin{verbatim}
Theorem omegabar_del_vt :
  vertex_transitiveb (T : diGraphType) ->
  forall u v : T, ω̄(del_tournament u) = ω̄(del_tournament v).
\end{verbatim}
\href{https://github.com/LLM4Rocq/digraph-theory/blob/main/theories/invariants_advanced/transitive.v\#L47}{Source: \code{theories/invariants\_advanced/transitive.v:47}}


\end{theorem}

\section{Domination bounds the clique number}

\begin{theorem}[\code{domnum\_le\_omegabar}]
  \label{res:domnum_le_omegabar}
  \uses{def:arc, def:backedge, def:domination, def:ltp, def:omegabar, def:tournament}
  \rocq{Digraph.invariants.domination.domnum_le_omegabar}
  \rocqok

\begin{verbatim}
Theorem domnum_le_omegabar : domnum <= ω̄(T).
\end{verbatim}
\href{https://github.com/LLM4Rocq/digraph-theory/blob/main/theories/invariants/domination.v\#L114}{Source: \code{theories/invariants/domination.v:114}}


\paragraph{Decoded.} $\mathrm{dom}(T) \le \bar\omega(T)$ for every tournament: a minimum
backedge-clique under an optimal order dominates
(\ref{def:domination}). This is the standard lower-bound tool for
$\bar\omega$.

\end{theorem}

\section{$C_3$ is the unique 2-critical tournament}

\begin{theorem}[\code{kcritical2\_uniq}]
  \label{res:kcritical2_uniq}
  \uses{def:C3, def:arc, def:dgiso}
  \rocq{Digraph.invariants.critical.kcritical2_uniq}
  \rocqok

\begin{verbatim}
Theorem kcritical2_uniq : dgiso (C3 : diGraphType) T.
\end{verbatim}
\href{https://github.com/LLM4Rocq/digraph-theory/blob/main/theories/invariants/critical.v\#L124}{Source: \code{theories/invariants/critical.v:124}}


\paragraph{Decoded.} Every $2$-$\bar\omega$-critical tournament is isomorphic
(\ref{def:dgiso}) to the directed triangle (\ref{def:C3}): the
base of the criticality hierarchy is unique.

\end{theorem}

\section{Proper subtournaments of critical tournaments}

\begin{theorem}[\code{kcritical\_proper\_sub}]
  \label{res:kcritical_proper_sub}
  \uses{def:arc, def:backedge, def:kcritical, def:ltp, def:omegabar, def:subdel, def:tournament}
  \rocq{Digraph.invariants.critical.kcritical_proper_sub}
  \rocqok

\begin{verbatim}
Lemma kcritical_proper_sub (k : nat) (T : tournament) (S : {set T}) :
  kcritical k T -> S != [set: T] -> (ω̄(sub_tournament S) <= k.-1)%N.
\end{verbatim}
\href{https://github.com/LLM4Rocq/digraph-theory/blob/main/theories/invariants/critical.v\#L46}{Source: \code{theories/invariants/critical.v:46}}


\paragraph{Decoded.} If $T$ is $k$-$\bar\omega$-critical and $S \subsetneq V(T)$, then the
induced subtournament on $S$ has $\bar\omega \le k - 1$: criticality
leaves no proper witness inside. This single lemma is what refutes
Question 5.9 at every $k$ (\ref{ch:unified}).

\end{theorem}

\section{$\bar\omega = 1$ exactly for transitive tournaments}

\begin{theorem}[\code{omegabar\_transb}]
  \label{res:omegabar_transb}
  \uses{def:arc, def:backedge, def:ltp, def:omegabar, def:tournament, def:transb}
  \rocq{Digraph.invariants.omegabar.omegabar_transb}
  \rocqok

\begin{verbatim}
Lemma omegabar_transb : 0 < #|T| -> (ω̄(T) == 1) = transb T.
\end{verbatim}
\href{https://github.com/LLM4Rocq/digraph-theory/blob/main/theories/invariants/omegabar.v\#L74}{Source: \code{theories/invariants/omegabar.v:74}}


\paragraph{Decoded.} For a nonempty tournament, $\bar\omega(T) = 1$ iff $T$ is transitive
(\ref{def:transb}) --- the sanity anchor for reading $\bar\omega$ as
a distance from transitivity. (The statement is a Boolean equation
between $\bar\omega(T) = 1$ and \code{transb T},
both Booleans.)

\end{theorem}

